\documentclass[12pt,letterpaper]{article}
\usepackage{graphicx,textcomp}
\usepackage{natbib}
\usepackage{setspace}
\usepackage{fullpage}
\usepackage{color}
\usepackage[reqno]{amsmath}
\usepackage{amsthm}
\usepackage{fancyvrb}
\usepackage{amssymb,enumerate}
\usepackage[all]{xy}
\usepackage{endnotes}
\usepackage{lscape}
\newtheorem{com}{Comment}
\usepackage{float}
\usepackage{hyperref}
\newtheorem{lem} {Lemma}
\newtheorem{prop}{Proposition}
\newtheorem{thm}{Theorem}
\newtheorem{defn}{Definition}
\newtheorem{cor}{Corollary}
\newtheorem{obs}{Observation}
\usepackage[compact]{titlesec}
\usepackage{dcolumn}
\usepackage{tikz}
\usetikzlibrary{arrows}
\usepackage{multirow}
\usepackage{xcolor}
\newcolumntype{.}{D{.}{.}{-1}}
\newcolumntype{d}[1]{D{.}{.}{#1}}
\definecolor{light-gray}{gray}{0.65}
\usepackage{url}
\usepackage{listings}
\usepackage{color}
\usepackage[parfill]{parskip}

\definecolor{codegreen}{rgb}{0,0.6,0}
\definecolor{codegray}{rgb}{0.5,0.5,0.5}
\definecolor{codepurple}{rgb}{0.58,0,0.82}
\definecolor{backcolour}{rgb}{0.95,0.95,0.92}

\lstdefinestyle{mystyle}{
	backgroundcolor=\color{backcolour},   
	commentstyle=\color{codegreen},
	keywordstyle=\color{magenta},
	numberstyle=\tiny\color{codegray},
	stringstyle=\color{codepurple},
	basicstyle=\footnotesize,
	breakatwhitespace=false,         
	breaklines=true,                 
	captionpos=b,                    
	keepspaces=true,                 
	numbers=left,                    
	numbersep=5pt,                  
	showspaces=false,                
	showstringspaces=false,
	showtabs=false,                  
	tabsize=2
}
\lstset{style=mystyle}
\newcommand{\Sref}[1]{Section~\ref{#1}}
\newtheorem{hyp}{Hypothesis}

\title{Problem Set 2}
\date{Due: February 18, 2026}
\author{Applied Stats II}


\begin{document}
	\maketitle
	\section*{Instructions}
	\begin{itemize}
		\item Please show your work! You may lose points by simply writing in the answer. If the problem requires you to execute commands in \texttt{R}, please include the code you used to get your answers. Please also include the \texttt{.R} file that contains your code. If you are not sure if work needs to be shown for a particular problem, please ask.
		\item Your homework should be submitted electronically on GitHub in \texttt{.pdf} form.
		\item This problem set is due before 23:59 on Wednesday February 18, 2026. No late assignments will be accepted.

	\end{itemize}

	\vspace{.25cm}

\noindent We're interested in what types of international environmental agreements or policies people support (\href{https://www.pnas.org/content/110/34/13763}{Bechtel and Scheve 2013)}. So, we asked 8,500 individuals whether they support a given policy, and for each participant, we vary the (1) number of countries that participate in the international agreement and (2) sanctions for not following the agreement. \\

\noindent Load in the data labeled \texttt{climateSupport.RData} on GitHub, which contains an observational study of 8,500 observations.

\begin{itemize}
	\item
	Response variable: 
	\begin{itemize}
		\item \texttt{choice}: 1 if the individual agreed with the policy; 0 if the individual did not support the policy
	\end{itemize}
	\item
	Explanatory variables: 
	\begin{itemize}
		\item
		\texttt{countries}: Number of participating countries [20 of 192; 80 of 192; 160 of 192]
		\item
		\texttt{sanctions}: Sanctions for missing emission reduction targets [None, 5\%, 15\%, and 20\% of the monthly household costs given 2\% GDP growth]
		
	\end{itemize}
	
\end{itemize}

\newpage
\noindent Please answer the following questions:

\begin{enumerate}
	\item
	Remember, we are interested in predicting the likelihood of an individual supporting a policy based on the number of countries participating and the possible sanctions for non-compliance.
	\begin{enumerate}
		\item [] Fit an additive model. Provide the summary output, the global null hypothesis, and $p$-value. Please describe the results and provide a conclusion.

	\end{enumerate}
	
	\item
	If any of the explanatory variables are statistically significant in this model, then:
	\begin{enumerate}
		\item
		For the policy in which nearly all countries participate [160 of 192], how does increasing sanctions from 5\% to 15\% change the odds that an individual will support the policy? (Interpretation of a coefficient)
		\item
		For the policy in which very few countries participate [20 of 192], how does increasing sanctions from 5\% to 15\% change the odds that an individual will support the policy? (Interpretation of a coefficient)
		\item
		What is the estimated probability that an individual will support a policy if there are 80 of 192 countries participating with no sanctions? 
	
	\end{enumerate}
	\item
	Would the answers to 2a and 2b potentially change if we included an interaction term in this model? Why? 
	\begin{itemize}
		\item Perform a test to see if including an interaction is appropriate.
	\end{itemize}
	\end{enumerate}

\newpage

\section{Question 1}

\subsection{Recoding}
Before performing any statistical analysis I check the data. I find that there are no NA values, however, the structure of the countries and sanctions variables as ordered factors would fit polynomial functions when input into the model. Thus I recode the variables to be regular factors with three and four levels respectively. 

\subsection{Model 1: Additive}
I fit an additive, logistic regression model  because the outcome variable is a dummy variable (a binary choice between support/non-support). Is it fair to say this about the regression output?

The output in \textbf{Table 1} shows how the log odds are effected by each predictor variable (when holding the other constant). It shows that the log odds have a greater increase the more countries involved, holding sanctions constant. However, compared to no sanctions and holding countries constant, the model shows that the log odds increase with a 5\% sanction, but further sanctions of 15\% or 20\% are associated with a decrease in the log odds of support.
It shows that there is a statistically significant relationship at all levels of both predictor variables, when holding the other constant. Thus we can reject the null hypothesis that the level of sanctions/ the number of countries involved in a climate initiative have no relationship with the probability of someone supporting a climate initiative.

\subsection{Global Null}
The global null is when all coefficients are 0, that all predictors have no relationship to the outcome variable. In this case H0 is that support for a climate initiatives have no relationship with the number of countries participating and the level of sanctions for missing emission targets. Ha, which is tested by model 1, is the alternative hypothesis, arguing that either predictor is associated with an change in support for a climate initiative.

To test the global null I run a likelihood ratio test to whether model 1 fits the data better. Effectively, this test shows whether the more complex model fits the data better than the model without any parameters. The p value of the LRT is statistically significant at \textit{\textbf{2.2e-16}}  (approximately 0.00000000000000022), while the deviance is \textbf{\textit{215.15}}. This suggests that model1, the model including the predictors fits the data better, indicating support for a climate initiative is associated with at least one of the predictors.

\subsection{R-code for models and LRT}
\begin{verbatim}
# Additive logistic model
model1 <- glm(choice ~ countries + sanctions,
              data = climateSupport,
              family = binomial)
              
# Global null model
model0 <- glm(choice ~ 1,
              data = climateSupport,
              family = binomial)
              
#Likelihood ratio test
anova(model0, model1, test = "LRT")
\end{verbatim}

\subsection{Additive logistic model output}
\begin{table}[!htbp] \centering 
 \caption{Sanctions and country participation on support for climate initiatives} 
  \label{} 
\begin{tabular}{@{\extracolsep{5pt}}lD{.}{.}{-3} } 
\\[-1.8ex]\hline 
\hline \\[-1.8ex] 
 & \multicolumn{1}{c}{\textit{Dependent variable:}} \\ 
\cline{2-2} 
\\[-1.8ex] & \multicolumn{1}{c}{choice} \\ 
\hline \\[-1.8ex] 
 Constant & -0.273^{***} \\ 
  & p = 0.00000 \\ 
  & \\ 
 countries 80 of 192 & 0.336^{***} \\ 
  & p = 0.000 \\ 
  & \\ 
 countries 160 of 192 & 0.648^{***} \\ 
  & p = 0.000 \\ 
  & \\ 
 sanctions 5\% & 0.192^{***} \\ 
  & p = 0.003 \\ 
  & \\ 
 sanctions 15\% & -0.133^{**} \\ 
  & p = 0.032 \\ 
  & \\ 
 sanctions 20\% & -0.304^{***} \\ 
  & p = 0.00001 \\ 
  & \\ 
\hline \\[-1.8ex] 
Observations & \multicolumn{1}{c}{8,500} \\ 
Log Likelihood & \multicolumn{1}{c}{-5,784.130} \\ 
Akaike Inf. Crit. & \multicolumn{1}{c}{11,580.260} \\ 
\hline 
\hline \\[-1.8ex] 
\textit{Note:}  & \multicolumn{1}{r}{$^{*}$p$<$0.1; $^{**}$p$<$0.05; $^{***}$p$<$0.01} \\ 
\end{tabular} 
\end{table} 
\newpage

\section{Question 2}
All explanatory variables in \textbf{Table 1} therefore all questions can be answered. Crucially however, the answer to (a) and (b) are the same, as the additive model holds countries constant at sanctions 5\% and 15\%.  This means that the change in odds that an individual will support a policy when increasing sanctions from 5\% to 15\% is the same regardless of the number of countries.

The change in odds can thus be calculated using the coefficients from \textbf{Table 1}. This is illustrated in the R-code below. The odds ratio is 0.7224479 (rounded to 0.722), a decrease in odds of about 27.8\%. This means, on average, increasing sanctions from 5\% to 15\% is associated with a 27.8\% lower odds of support, holding number of countries constant.

\textbf{R-code}
\begin{verbatim}
# check coefficients
summary(model1)
# 15% is -0.13325, 5% is 0.19186 - simple maths finds the difference in log odds
-0.13325 - 0.19186
# answer is a log odds of -0.32511
# finding the odds ratio
exp(-0.32511)
# odds ratio is 0.7224479
\end{verbatim}

\textbf{(c)} differs to the extent that it cannot be answered via \textbf{Table 1} alone. It requires using the predict function to find the probability based on the previous model. Using this function the difference in log odds (moving from 20 to 80 countries) is 0.06369783 and the probability of support is 0.5159191 (about 51.59\%). This means, that the estimated probability that an individual will support a policy if there are 80 countries participating and there are no sanctions for not meeting emissions targets is approximately 51.59\%.

\textbf{R-code}
\begin{verbatim}
# for log odds
predict(model1,
        newdata = data.frame(countries="80 of 192", sanctions="None"),
        type="link")
# answer = 0.06369783 change in log odds

# for probability
predict(model1,
        newdata = data.frame(countries="80 of 192", sanctions="None"),
        type="response")
# answer = 0.5159191 or a 51.59% probability of support.
\end{verbatim}

\section{Question 3}
Theoretically the answers could change if the effect of sanctions were to vary by number of participant countries. To test this, I fit a model with an interaction between the sanctions and countries variable. This is \textbf{model 2} and is compared to \textbf{model 1} in the below \textbf{Table 2}. Model2 shows minimal effects on the log odds of the number of participant countries on support for the initiative when accounting for the interaction effect between countries and sanctions. At a glance, this would suggest that there would not be a significant divergence from 2a and 2b.

To further test this, I will perform a likelihood ratio test, to test whether including the interaction improves the model fit. Acknowledging that we seek a parsimonious model, the decrease in deviance of 6.29, while increasing the degrees of freedom by 6, suggests that the interaction model does not greatly improve the model fit, while making it more complex. Additionally, the p-value of 0.391 is not significant, thus we fail to reject the null hypothesis -- that interaction effect has no effect on support. This suggests that \textbf{model 1}, the additive model is prefered, and that there is no evidence that 2a and 2b would change.

\begin{table}[!htbp] \centering 
  \caption{Sanctions and country participation on support for climate initiatives: models compared} 
  \label{} 
\begin{tabular}{@{\extracolsep{5pt}}lD{.}{.}{-3} D{.}{.}{-3} } 
\\[-1.8ex]\hline 
\hline \\[-1.8ex] 
 & \multicolumn{2}{c}{\textit{Dependent variable:}} \\ 
\cline{2-3} 
\\[-1.8ex] & \multicolumn{2}{c}{choice} \\ 
\\[-1.8ex] & \multicolumn{1}{c}{(1)} & \multicolumn{1}{c}{(2)}\\ 
\hline \\[-1.8ex] 
 Constant & -0.273^{***} & -0.275^{***} \\ 
  & (0.054) & (0.075) \\ 
  countries80 of 192 & 0.336^{***} & 0.376^{***} \\ 
  & (0.054) & (0.106) \\ 
  countries160 of 192 & 0.648^{***} & 0.613^{***} \\ 
  & (0.054) & (0.108) \\ 
  sanctions5\% & 0.192^{***} & 0.122 \\ 
  & (0.062) & (0.105) \\ 
  sanctions15\% & -0.133^{**} & -0.097 \\ 
  & (0.062) & (0.108) \\ 
  sanctions20\% & -0.304^{***} & -0.253^{**} \\ 
  & (0.062) & (0.108) \\ 
  countries80 of 192:sanctions5\% &  & 0.095 \\ 
  &  & (0.152) \\ 
  countries160 of 192:sanctions5\% &  & 0.130 \\ 
  &  & (0.151) \\ 
  countries80 of 192:sanctions15\% &  & -0.052 \\ 
  &  & (0.152) \\ 
  countries160 of 192:sanctions15\% &  & -0.052 \\ 
  &  & (0.153) \\ 
  countries80 of 192:sanctions20\% &  & -0.197 \\ 
  &  & (0.151) \\ 
  countries160 of 192:sanctions20\% &  & 0.057 \\ 
  &  & (0.154) \\ 
 \hline \\[-1.8ex] 
Observations & \multicolumn{1}{c}{8,500} & \multicolumn{1}{c}{8,500} \\ 
Log Likelihood & \multicolumn{1}{c}{-5,784.130} & \multicolumn{1}{c}{-5,780.983} \\ 
Akaike Inf. Crit. & \multicolumn{1}{c}{11,580.260} & \multicolumn{1}{c}{11,585.970} \\ 
\hline 
\hline \\[-1.8ex] 
\textit{Note:}  & \multicolumn{2}{r}{$^{*}$p$<$0.1; $^{**}$p$<$0.05; $^{***}$p$<$0.01} \\ 
\end{tabular} 
\end{table}

\textbf{R-code}
\begin{verbatim}
# model with interaction term 
model2 <- glm(choice ~ countries*sanctions,
              data = climateSupport,
              family = binomial)
# summarise model
summary(model2)

# run stargazer
stargazer(model1, model2,
          type = "latex",
          title="Sanctions and country participation on support for climate initiatives: models compared",
          align=TRUE,
          no.space = TRUE,
          intercept.top = TRUE,
          intercept.bottom = FALSE)

#LRT for both models - does model2 improve fit.
anova(model1, model2, test = "LRT")
\end{verbatim}

\end{document}
